
\section*{Alban au XV\textsuperscript{ème} siècle}

\lettrine{U}{ne} tradition prétend que c'est à ce moment qu'aurait été
planté le « hêtre de Margou » qu'on continue à appeler
« lou faous dé Roussignol ». Même si le cadastre de 1598 ne
fait pas état d'une famille Rossignol -- le propriétaire de
cette époque pouvant avoir changé de nom -- cette tradition
demeure valable tant qu'elle n'est pas infirmée par quelque
témoignage contraire. Cet arbre illustre continue d'ailleurs
à figurer sur nos cartes comme point géodésique. Et puisque
nous parlons environnement vert, peut-être est-il légitime
de penser que remonte jusqu'à Henri IV l'arbre quelque peu
fatigué entouré d'un banc de pierre aussi vétuste que lui,
qui orne la petite place de Notre Dame d'Ourtiquet. Ce
pourrait bien être un de ces ormeaux dits de « Sully ».

\lettrine{D}{e} ce temps bienheureux reste un parchemin
(versé aux archives départementales) de 1547, ayant
appartenu à la famille de Louis Salluste, le sympathique
garde champêtre d'Alban. Par sa maman, ses ancêtres furent
les « Pujol de la Tapie » (paroisse de Notre Dame d'Alban,
mais commune de Massals). Cette famille fera
particulièrement parler d'elle à l'occasion des guerres de
religion. L'un de ces « Pujol de la Taie\footnote{Tapie?} » est en 1547 curé
de Cambon près d'Albi. L'âge et la maladie lui laissant
deviner que sa fin ne saurait trop tarder, c'est par devant
Me Duin, notaire à Albi qu'il exprime ses dernières volontés.
A l'origine, ce document était constitué par sept peaux
d'agneau cousues à la suite l'une de l'autre : ce qui faisait
une bande d'environ 3 mètres à 3 mètres cinquante sur
0,60 de largeur. L'original fut par suite mutilé, il en reste
environ la moitié. Le texte était rédigé en occitan
(=patois). Une copie du texte intégral demeure cependant,
et peut être consultée aux archives départementales à
Albi. Le testateur y règle, non seulement l'attribution de
ses biens : maison, forge, chais, vignobles... etc... mais
aussi ses obsèques et l'emplacement de sa tombe dans
l'église de Cambon d'Albi. Jusqu'à ce jour, l'auteur de ces
